\minitopic{研读实验室：同一证据为何得到不同判断}\index{语境主义}\index{实践侵入}\index{断言规范}\index{行动规范} 银行案例之所以重要，不是它证明利害一定改变知识，而是它把三个变量放在一起：说话语境、主体风险与行动要求。若周六是否营业只影响一次散步，主体凭上周营业经历说“我知道”；若房款逾期会造成重大损失，配偶要求再次确认。直觉变化可能来自“知道”的使用标准提高，也可能来自主体需要更强证据，或只是行动值得复核。理论比较必须设法把三者分开。 最小对比法要求一次改变一项。第一组保持主体证据和利害不变，只改变谈话中是否提到特殊停业可能，用来检验语境主义。第二组保持说话者用词与听者标准不变，只改变主体若错将承担的损失，用来检验实践侵入。第三组明确说主体仍知道，却问他是否应在行动前花一分钟查官网，用来检验“高风险只增加行动审慎”的解释。若三组判断都变化，可能不止一个机制在工作。 还要控制情绪与道德责备。高风险故事常写得更生动，使读者因同情而降低知识归属；主体若有轻松核验机会却不使用，也容易被视为粗心。实验上应设置同样的核验成本、同样的主体谨慎度，只改变损失。哲学上则应分别询问：“他知道吗”“他相信得合理吗”“他应该复核吗”“他若不复核是否可责”。

\minitopic{语境主义的语言检验}语境主义是关于知识归属句的语义理论，不能仅凭人们在不同故事中给出不同判断而成立。需要考察语言行为。例如，低标准语境中的甲说“李知道银行营业”，高标准语境中的乙随后说“甲说的是真的，但在我们这里李不知道”，这段跨语境报告是否自然？若同一归属在语境切换后可保持真值，语义变化解释会受压力。 另一个检验是显式标准。说“按日常意义他知道，按绝对确定意义他不知道”很自然，但这也可能只是“知道”有松紧用法，并不证明普通词具有系统的语境敏感语义。语境主义需说明标准由哪些会话规则改变：显著替代、目的、听者问题，还是说话者意图。单说“语境决定”没有预测力。 怀疑论情景的提出尤其复杂。若只要有人说“也许银行网站被黑”，便能取消一切知识，恶意参与者可无限抬高标准。Lewis 的相关替代思路试图用忽略规则限制何种可能不可被正当忽略\parencite{lewis1996}。但规则本身也需说明：当前证据使可能相关，还是会话注意使它相关？纯注意若足够，知识归属会过度易变。 语境主义的优势是保留日常可错知识与哲学高标准直觉。它的代价不是“相对主义”这个标签，而是必须给出足够具体的语义证据，并解释主体自己用第一人称“我知道”时标准如何确定。语义理论若只复述案例直觉，就未完成解释。

\minitopic{实践侵入的论证结构}实践侵入主张，传统的真值、信念和证据因素完全相同，主体是否知道仍可能因实际利害而不同\parencite{stanley2005}。一条典型论证从知识与合理行动连接出发：（1）若主体知道$p$，就可把$p$作为行动前提；（2）高风险情形中主体不可仅凭现有根据把$p$作为行动前提；所以（3）他不知道$p$。争点集中在第一前提是否无例外，以及第二前提是否真的说明知识不足，而非行动需要冗余保障。 消防员可能知道楼内无人，却因错误后果极重仍执行二次搜索。复核不一定表明先前无知识；它可能是安全协议对低概率灾难的成本敏感响应。相反，若主体的注意和取证行为会受利害影响，高风险可能改变实际证据：他意识到过去忽略的节假日、想起信息过期，或获得关于自己粗心的高阶证据。此时知识变化不需要实践因素直接进入知识条件。 实践侵入的支持者可回应，知识不仅是静态证据地位，也标志主体何时可以结束探究。若在重大决定中仍必须继续调查，称其已经知道似乎失去行动功能。反对者则区分“可合理结束求真探究”与“在多种价值约束下可停止风险管理”。双方真正分歧是知识与行动理由的规范联系，不只是银行故事的语感。

\minitopic{知识、断言与限定表达}知识规范常被表述为：直接断言$p$只有在知道$p$时才适当。它解释了质问“你怎么知道？”为何能挑战断言，也解释说话者发现自己并不知道后为何应撤回。但断言有多种目的：课堂上提出待检验假说、紧急情况下传递高概率警报、法庭上宣誓陈述事实，所需地位未必相同。 语言有丰富的校准工具。“证据显示”“我有七成把握”“截至目前没有发现”“若记录无误则”不是较弱写作，而是在标明证据类型和条件。把这些句子都当作断言$p$的失败版本，会损害科学交流。知识规范最适合约束无保留、以信息传递为目的的事实断言，而非所有言语行为。 说话者也可知道$p$却不应断言。例如信息属于隐私，公开会造成伤害；又如听者会系统误解缩略说法。这里道德或交流理由击败断言许可，但不击败知识。由此可见，“知道”即使是断言的认识规范，也不是断言的全部规范。 反方向上，不知道$p$有时仍应发出警告。毒气传感器尚未完成确认，却显示高风险异常；值班员说“可能泄漏，请撤离”是合适的。句子内容包含模态限定，其目标命题并非无保留的“已泄漏”。行动与语言可在不虚构知识的前提下响应不确定性。

\minitopic{比较视角：因时、因位与相对主义的区别}早期儒家文本经常把言说与行动放在角色、时机和后果中评价。《论语》中“知之为知之，不知为不知”的熟语强调如实校准自我陈述，而大量问答又因提问者与情境不同给出不同劝诫\parencite[2.17]{analects2003}。这可以与当代语境敏感问题对话，但不能直接推成一套关于“知道”语义的语境主义。 “因时制宜”至少有三种读法。第一，事实不变而行动要求随角色变化；第二，相同证据对不同任务是否足够不同；第三，命题真值本身相对角色变化。前两项不蕴涵第三项。医生和患者承担不同说明义务，不等于影像是否有病灶因身份而真或假；法官要求排除合理怀疑，也不等于被告是否实施行为由法庭标准创造。 比较的价值在于打破一个现代假设：知识评价可以完全脱离说话、行动与角色责任。与此同时，文本边界防止把每种情境化判断都叫“相对主义”。较稳妥的结论是，不同传统都能提示认识规范具有实践嵌入性；它们如何解释知识条件，仍需分别重构。

\minitopic{综合案例：灾害预警的发布门槛}气象部门依据多个模型判断台风将转向人口密集区，综合概率为70\%。内部专家是否“知道”台风会转向？若知识要求命题为真，在事件发生前的概率证据是否足够仍有争议；但部门必须决定是否发布预警。漏报造成生命损失，误报造成撤离成本和长期信任下降。 语境主义会研究专家会议、媒体发布和日常谈话中知识归属标准是否不同；实践侵入会问巨大损失是否使同一证据不足以知道；不变主义者可保持知识条件不变，把差异放在行动阈值。决策理论则要求把概率、损失和可逆性纳入行动，不必先裁定知识归属。 好的发布可以说：“现有模型有七成指向转向；主要不确定性来自海温路径；下一次更新在两小时后；鉴于撤离需要六小时，现在发布二级预警。”这段话把证据、未知和行动理由分开。若后来台风未转向，不能仅凭结果判定预警不理性；应检查概率是否校准、模型是否独立、当时决策阈值是否合理。 若部门内部已有新观测把概率降至20\%却因程序延迟未进入报告，知识与责任评价改变：对使用旧报告的公众，证言依赖可能仍合理；对知晓更新却沉默的管理者，高阶证据构成击败。语境、主体信息与制度角色在同一案例中发生作用，不能以单一“高风险”解释全部。

\begin{tcolorbox}[breakable,colback=white,colframe=blue!30!black,title={本章研读任务},fonttitle=\bfseries]
设计三版同证据案例：只改变会话中显著替代，只改变主体损失，只改变核验成本。分别记录知识、合理信念、行动和责任判断；再用语境主义、实践侵入与行动审慎解释，指出哪一版最能区分理论。
\end{tcolorbox}
